% ============ 第六章 §1 Hasse 类数公式 (原书 pp.206-216) ============
\section{Hasse 类数公式}\label{sec:ch6-s1}

对于 Abel 数域 $K$, $\zeta_K(s)$ 应当有好的表达式, 这一点是不应当感到很奇怪的. 因为根据 Hronecker-Weber 定理, 每个 Abel 数域均是分圆域的子域, 而分圆域中有很简单的素理想分解规则. 设 $\Q(\zeta_m)(m\not\equiv 2\ (\mathrm{mod}\ 4))$ 是包含 Abel 数域 $K$ 的最小分圆域, 即 $m=\operatorname{cond}(K)$（域 $K$ 的导子）. 我们有如下的伽罗瓦对应:
\[
  \begin{array}{cc}
    \Q(\zeta_m) & G=\Gal(\Q(\zeta_m)/\Q)\\
    \mid & \mid\\
    K & H\\
    \mid & \mid\\
    \Q & \{1\}
  \end{array}
\]
由于
\[
  G\cong(\Z/m\Z)^\times,\qquad\sigma_a\mapsto a\ (\mathrm{mod}\ m),\quad (a,m)=1,\quad \sigma_a(\zeta_m)=\zeta_m^a.
\]
从而 $H$ 同构于 $(\Z/m\Z)^\times$ 的一个子群（我们今后常常把 $H$ 和这个子群等同起来）, 而 $\Gal(K/\Q)$ 同构于商群 $(\Z/m\Z)^\times/H$（我们也常把这两者等同起来）. 有限 Abel 群 $\Gal(K/\Q)=(\Z/m\Z)^\times/H$ 的每个特征 $\chi$ 也叫作是域 $K$ 的特征, 而 $(\Z/m\Z)^\times/H$ 的特征群也叫作是 $K$ 的特征群, 表示成 $\widehat{K}$. 于是 $\widehat{K}$ 实际上是由全部在 $H$ 上平凡的模 $m$ 的 D-特征所构成的.

设 $\chi$ 为模 $m$ 的 D-特征, 则 $\chi$ 是由惟一决定的一个模 $m'$ 的本原 D-特征所诱导出来的, 其中 $m'=\operatorname{cond}(\chi)$, $m'\mid m$. 我们将诱导出 $\chi$ 的这个本原 D-特征表示成 $\chi^*$.

如果 $\Q\subseteq K\subseteq L\subseteq\Q(\zeta_m)$, 而 $E_K$ 和 $E_L$ 分别为 $G=\Gal(\Q(\zeta_m)/\Q)$ 对于 $K$ 和 $L$ 的固定子群, 则 $E_K\subseteq E_L$, 即伽罗瓦对应 $K\leftrightarrow E_K$ 是反序的（序指的是集合之间的包含关系）. 由于 $\Gal(K/\Q)\simeq G/E_K$, $\Gal(L/\Q)\cong G/E_L$, 可知它们的特征群之间有关系 $\widehat{K}\subseteq\widehat{L}$. 所以对应 $K\mapsto\widehat{K}$ 是保序的.

群 $G$ 同构于 $(\Z/m\Z)^\times$. 若 $m=p_1^{e_1}\cdots p_r^{e_r}$ 是 $m$ 的素因子分解式, 则有同构
\[
  (\Z/m\Z)^\times\cong(\Z/p_1^{e_1}\Z)^\times\times\cdots\times(\Z/p_r^{e_r}\Z)^\times\qquad(\text{直积}),
\]
所以每个模 $m$ 的 D-特征惟一分解为 $\chi=\chi_{p_1}\cdots\chi_{p_r}$, 其中 $\chi_{p_i}$ 是模 $p_i^{e_i}$ 的特征（$1\leqslant i\leqslant r$）. 易知 $f_\chi=f_{\chi_{p_1}}\cdots f_{\chi_{p_r}}$, 这里 $f_\chi$ 表示特征 $\chi$ 的导子. 如果 $p$ 不为 $p_i(1\leqslant i\leqslant r)$, 则 $\chi_p$ 是主特征. 现在对 Abel 数域 $K$ 和每个素数 $p$, 定义如下一些特征集合:
\[
  \widehat{K}_p=\{\chi_p:\chi\in\widehat{K}\},\qquad
  Y=Y(K,p)=\{\chi\in\widehat{K}:\chi^*(p)\neq 0\},\qquad
  Z=\{\chi\in Y:\chi^*(p)=1\}.
\]
易知它们均形成特征群, 并且 $Z\subseteq Y\subseteq\widehat{K}$. 下面引理表明: 这些特征群可以刻画 $p$ 在 $K$ 中的分解型式.

\begin{lemma}\label{lem:ch6-1}
以 $e,f,g$ 分别表示 $p$ 在 Abel 数域 $K(\subseteq\Q(\zeta_m))$ 中的分歧指数, 剩余类域次数和分解次数, 则
\begin{enumerate}[(1)]
  \item $e=|\widehat{K}_p|$, 并且 $p$ 在 $K$ 中分歧当且仅当存在 $\chi\in\widehat{K}$, 使得 $\chi^*(p)=0$（即 $p\mid f_\chi$）.
  \item $Y$ 和 $Z$ 分别是 $p$ 对 $K$ 的惯性域和分解域的特征群. 于是
    \[
      e=[\widehat{K}:Y],\qquad f=[Y:Z],\qquad g=|Z|.
    \]
\end{enumerate}
\end{lemma}

\begin{proof}
(1) 设 $m=p^ln$, $(p,n)=1$, $l\geqslant 0$, 则有如下的域扩张图表:
\[
  \begin{array}{cc}
    \Q(\zeta_m)=\Q(\zeta_{p^l})\Q(\zeta_n) & (\Z/n\Z)^\times\times\widehat{K}_p\\
    \mid & \mid\\
    L=\Q(\zeta_n)K & \widehat{L}\\
    \mid & \mid\\
    K & \widehat{K}\\
    \mid & \mid\\
    \Q & \{1\}
  \end{array}
\]
图中每个域的特征群均看成是 $\Q(\zeta_m)$ 的特征群（即模 $m$ 的 D-特征群）的子群. 由于 $L$ 是 $\Q(\zeta_n)$ 和 $K$ 的合成, 而域到此域特征群之间的一一对应是保序的, 所以 $\widehat{L}$ 是由 $\Q(\zeta_n)$ 和 $\widehat{K}$ 生成的特征群. 于是
\[
  \widehat{L}=\langle(\Z/n\Z)^\times,\widehat{K}\rangle=(\Z/n\Z)^\times\ \text{和}\ \widehat{K}_p\ \text{的直积}.
\]
以 $F$ 表示对应于 $\widehat{K}_p$ 的域, 即 $\widehat{F}=\widehat{K}_p$, 则 $L=\Q(\zeta_n)F$, 并且 $F\subseteq\Q(\zeta_{p^l})$. 对于 $\Q(\zeta_m)$ 的每个子域 $M$, 以 $e(M)$ 表示 $p$ 在 $M$ 中的分歧指数, 则
\begin{align*}
  e=e(K)&=e(L)\qquad(\text{因为 } e(\Q(\zeta_n))=1\ \text{而 } L=\Q(\zeta_n)K)\\
  &=e(F)\qquad(\text{因为 } L=\Q(\zeta_n)F\ \text{而 } e(\Q(\zeta_n))=1)\\
  &=[F:\Q]\qquad(\text{因为 } F\subseteq\Q(\zeta_{p^l}),\ \text{从而 } p\ \text{在 } F\ \text{中完全分歧})\\
  &=|\widehat{F}|=|\widehat{K}_p|.
\end{align*}
进而, $p$ 在 $K$ 中分歧 $\Leftrightarrow|\widehat{K}_p|=e\geqslant 2\Leftrightarrow$ 存在 $\chi\in\widehat{K}$, 使得 $\chi_p\neq\chi_0\Leftrightarrow$ 存在 $\chi\in\widehat{K}$, 使得 $\chi^*(p)=0$. 这就证明了 (1).

(2) 设 $M$ 为 $K$ 的任一子域. 由 (1) 可知
\[
  p\ \text{在 } M\ \text{中不分歧}\Leftrightarrow\ \text{对每个 } \chi\in\widehat{M},\ \text{均有 } \chi^*(p)\neq 0
  \Leftrightarrow\widehat{M}\subseteq Y.
\]
这就表明 $Y$ 对应于 $p$ 在 $K$ 中的最大不分歧子域, 即对应于 $p$ 在 $K$ 中的惯性域 $K_I$. 所以
\[
  Y=\widehat{K}_I,\qquad e=[K:K_I]=[\widehat{K}:\widehat{K}_I]=[\widehat{K}:Y].
\]
进而, 对每个 $\chi\in Y$, 由 $\chi^*(p)\neq 0$ 可知 $f_\chi\mid n$. 所以 $\chi$ 是由一个确定的模 $n$ 特征 $\chi'$ 所诱导的, 并且 $(\chi')^*=\chi^*$. 另一方面, $\Q(\zeta_n)$ 是 $p$ 在 $\Q(\zeta_m)$ 中的最大不分歧子域, 所以 $K_I\subseteq\Q(\zeta_n)$. 以 $K_D$ 表示 $p$ 在 $K$ 中的分解域, 于是有如下的图表:
\[
  \begin{array}{cc}
    K & \Gal(K/\Q)\\
    \mid & \mid\\
    K_I & \Gal(K_I/\Q)\\
    \mid & \mid\\
    K_D & \Gal(K_D/\Q)\\
    \mid & \mid\\
    \Q & \{1\}
  \end{array}
\]
其中 $I$ 和 $D$ 分别为对应于 $K_I$ 和 $K_D$ 的 $(\Z/n\Z)^\times$ 的子群. 而 $K_I$ 看作 $\Q(\zeta_n)$ 的子域时, 其特征群为 $Y'=\{\chi':\chi\in Y\}$. 由于
\[
  \Gal(K_D/\Q)=(\Z/n\Z)^\times/D\cong\frac{(\Z/n\Z)^\times/I}{D/I},
\]
从而对每个 $\chi'\in Y'=((\Z/n\Z)^\times/I)^{\wedge}$,
\begin{align*}
  \chi'\ \text{为 } \Gal(K_D/\Q)\ \text{的特征}&\Leftrightarrow\chi'(D/I)=1\\
  &\Leftrightarrow\chi'(p)=1\qquad(\text{因为 } D/I\ \text{是由 } \sigma_p\ \text{生成的循环群})\\
  &\Leftrightarrow\chi^*(p)=1\Leftrightarrow\chi\in Z.
\end{align*}
于是 $\widehat{K}_D=Z$, 所以 $f=[K_I:K_D]=[Y:Z]$ 并且 $g=|Z|$. 这就证明了引理.
\end{proof}

现在我们证明: 对每个 Abel 数域 $K$, Dedekind zeta 函数 $\zeta_K(s)$ 是一些对本原特征的 Dirichlet $L$-函数的乘积. 令 $m$ 为域 $K$ 的导子, 则 $K$ 的特征 $\chi$ 都可看成模 $m$ 的 D-特征, 从而 $L(s,\chi^*)$ 是 Dirichlet $L$-函数.

\begin{theorem}[\emph{Hasse}]\label{thm:6.1}
设 $K$ 是 Abel 数域, 则
\[
  \zeta_K(s)=\prod_{\chi\in\widehat{K}}L(s,\chi^*)=\zeta(s)\prod_{\substack{\chi\in\widehat{K}\\ \chi\neq\chi_0}}L(s,\chi^*),
\]
即 Dedekind zeta 函数 $\zeta_K(s)$ 是一些对本原 D-特征的 Dirichlet $L$-函数之乘积.
\end{theorem}

\begin{proof}
根据解析延拓原理, 我们只要对 $\operatorname{Re}(s)>1$ 的情形证明上式即可. 而在 $\operatorname{Re}(s)>1$ 时, 两边均有 Euler 乘积公式: $\zeta_K(s)=\prod_p\prod_{\mathfrak{p}\mid p}(1-N(\mathfrak{p})^{-s})^{-1}$, $L(s,\chi^*)=\prod_p(1-\chi^*(p)p^{-s})^{-1}$. 从而只需对每个有理素数 $p$ 证明
\begin{equation}\tag{1}
  \prod_{\mathfrak{p}\mid p}(1-N(\mathfrak{p})^{-s})=\prod_{\chi\in\widehat{K}}(1-\chi^*(p)p^{-s})
\end{equation}
即可. 设
\[
  pO_K=(\mathfrak{p}_1\cdots\mathfrak{p}_g)^e,\qquad N(\mathfrak{p}_i)=p^f\ (1\leqslant i\leqslant g),\qquad efg=n=[K:\Q],
\]
则 (1) 式左边为 $(1-p^{-fs})^g$. 另一方面, 利用引理 \ref{lem:ch6-1} 和其中的记号, 可知当 $\chi\notin Y$ 时, $\chi^*(p)=0$. 所以 (1) 式右边为 $\prod_{\chi\in Y}(1-\chi^*(p)p^{-s})$. 进而, $\chi\in Z\Leftrightarrow\chi^*(p)=1$. 并且 $Y/Z$ 是 $f$ 阶循环群, 从而 $Y$ 对 $Z$ 的每个陪集中的 $g$ 个特征 $\chi$, $\chi^*(p)$ 取值相同. 并且对不同陪集中的 $f$ 个特征 $\chi_1,\cdots,\chi_f$, 它们对应的本原特征 $\chi_1^*,\cdots,\chi_f^*$ 在 $p$ 的取值分别为 $1,\zeta_f,\cdots,\zeta_f^{f-1}$. 因此
\begin{align*}
  \prod_{\chi\in Y}(1-\chi^*(p)p^{-s})
  &=\prod_{\chi\in Y/Z}(1-\chi^*(p)p^{-s})^g\\
  &=\prod_{i=0}^{f-1}(1-\zeta_f^ip^{-s})^g=(1-p^{-fs})^g.
\end{align*}
这就证明了 (1) 式, 从而也完全证明了定理.
\end{proof}

定理 \ref{thm:4.14} 给出了 $\zeta_K(s)$ 与 $K$ 的理想类数 $h_K$ 的关系. 现在 $\zeta_K(s)$ 可以表成 Dirichlet $L$-函数的乘积, 所以 $h_K$ 可以用 Dirichlet $L$-函数来表达.

\begin{corollarynonum}
设 $K$ 为 Abel 数域, 则
\[
  \rho_Kh_K=\prod_{\substack{\chi\in\widehat{K}\\ \chi\neq\chi_0}}L(1,\chi^*),
\]
其中
\[
  \rho_K=\frac{2^{r_1}(2\pi)^{r_2}R_K}{w_K|d(K)|^{1/2}}.
\]
\end{corollarynonum}

\begin{proof}
定理 \ref{thm:4.14} 给出 $\rho_Kh_K=\lim\limits_{s\to1}\dfrac{\zeta_K(s)}{s-1}$. 而由定理 \ref{thm:6.1} 又知道
\[
  \zeta_K(s)=\prod_{\chi\in\widehat{K}}L(s,\chi^*)=\zeta(s)\prod_{\substack{\chi\in\widehat{K}\\ \chi\neq\chi_0}}L(s,\chi^*),
\]
其中 $\zeta(s)=L(s,\chi_0)$ 是黎曼 zeta 函数. 我们知道当 $\chi\neq\chi_0$ 时, $\chi^*\neq\chi_0$, 从而 $L(s,\chi^*)$ 在 $s=1$ 处正则. 而 $\lim\limits_{s\to1}\dfrac{\zeta(s)}{s-1}=1$. 因此
\[
  \lim_{s\to1}\frac{\zeta_K(s)}{s-1}=\prod_{\substack{\chi\in\widehat{K}\\ \chi\neq\chi_0}}L(1,\chi^*).
\]
这就证明了推论.
\end{proof}

\begin{sloppypar}
现在的问题是将 $L(1,\chi^*)$ 表示成尽可能初等和便于计算的形式. 注意对每个 Dirichlet 特征 $\chi$, $\chi(-1)^2=\chi((-1)^2)=\chi(1)=1$, 从而 $\chi(-1)=\pm 1$. 当 $\chi(-1)=1$ 时, $\chi$ 叫作偶特征, 而 $\chi(-1)=-1$ 时, $\chi$ 叫作奇特征. 今后将看出, 按照 $\chi$ 是奇特征还是偶特征, $L(1,\chi^*)$ 有不同形式的表达式. 首先我们指出特征 $\chi$ 的奇偶性的一个数论意义.
\end{sloppypar}

\begin{lemma}\label{lem:ch6-2}
(a) Abel 数域 $K$ 为实域 $\Leftrightarrow K$ 的特征均为偶特征.

(b) 若 $K$ 为虚 Abel 数域, $K_+$ 为 $K$ 的极大实子域, 则 $[K:K_+]=2$, 并且 $\widehat{K}_+$ 即为 $\widehat{K}$ 中全部偶特征所形成的子群.
\end{lemma}

\begin{proof}
(a) 假设 $\operatorname{cond}(K)=m$, 则有伽罗瓦对应
\[
  \begin{array}{cc}
    \Q(\zeta_m) & (\Z/m\Z)^\times\\
    \mid & \mid\\
    K & G\\
    \mid & \mid\\
    \Q & \{1\}
  \end{array}
\]
于是 $\Gal(K/\Q)=(\Z/m\Z)^\times/G$, 而 $(\Z/m\Z)^\times$ 中的 $-1\ (\mathrm{mod}\ m)$ 相当于复共轭自同构 $\sigma_{-1}(\zeta_m)=\zeta_m^{-1}=\overline{\zeta}_m$. 从而
\begin{align*}
  K\ \text{为实域}&\Leftrightarrow\sigma_{-1}\ \text{为 } K\ \text{中恒等自同构}\Leftrightarrow -1\in G\\
  &\Leftrightarrow -1\ \text{为 } (\Z/m\Z)^\times/G\ \text{中的单位元素}\\
  &\Leftrightarrow K\ \text{的每个特征（即 } (\Z/m\Z)^\times/G\ \text{的每个特征）均为偶特征}.
\end{align*}
(b) 若 $K$ 为虚 Abel 数域, 则 $K$ 中有奇特征, 于是 $\widehat{K}$ 中偶特征全体形成 $\widehat{K}$ 的指数为 $2$ 的子群. 此子群即是 $(\Z/m\Z)^\times/\langle -1,G\rangle$ 的特征群, 而 $\langle -1,G\rangle$ 的固定子域就是 $K$ 的极大实子域 $K_+$, 由此即得结论.
\end{proof}

由引理 \ref{lem:ch6-2} 我们有
\[
  \zeta_{K_+}(s)=\prod_{\chi\in\widehat{K}_+}L(s,\chi^*)=\prod_{\substack{\chi\in\widehat{K}\\ \chi(-1)=1}}L(s,\chi^*)=\zeta(s)\prod_{\substack{\chi\in\widehat{K}\\ \chi(-1)=1\\ \chi\neq\chi_0}}L(s,\chi^*).
\]
从而由定理 \ref{thm:6.1} 的系立刻得到

\begin{lemma}\label{lem:ch6-3}
(a) 设 $K$ 为 $n$ 次实 Abel 域, 则
\[
  R_Kh_K\cdot\frac{2^{n-1}}{|d(K)|^{1/2}}=\prod_{\substack{\chi\in\widehat{K}\\ \chi\neq\chi_0}}L(1,\chi^*).
\]
(b) 设 $K$ 为 $n$ 次虚 Abel 域, 则
\[
  R_Kh_K\frac{(2\pi)^{n/2}}{w_K|d(K)|^{1/2}}=\prod_{\substack{\chi\in\widehat{K}\\ \chi\neq\chi_0}}L(1,\chi^*),
\]
\[
  R_{K_+}h_{K_+}\cdot\frac{2^{n/2-1}}{|d(K_+)|^{1/2}}=\prod_{\substack{\chi_0\neq\chi\in\widehat{K}\\ \chi(-1)=1}}L(1,\chi^*).
\]
\end{lemma}

现在我们来计算 $L(1,\chi^*)$ 的值.

\begin{lemma}\label{lem:ch6-4}
设 $\chi$ 为模 $m$ 本原 D-特征, $m\geqslant 3$, 则
\[
  L(1,\chi)=\begin{cases}
    -\dfrac{2G(1,\chi)}{m}\sum\limits_{1\leqslant k<m/2}\overline{\chi}(k)\log\sin\dfrac{k\pi}{m}, & \text{当 } \chi(-1)=1\ \text{时};\\[1em]
    \dfrac{\pi iG(1,\chi)}{m^2}\sum\limits_{k=1}^{m-1}\overline{\chi}(k)k=\dfrac{\pi iG(1,\chi)}{m(\chi(2)-2)}\sum\limits_{1\leqslant k<m/2}\overline{\chi}(k), & \text{当 } \chi(-1)=-1\ \text{时},
  \end{cases}
\]
其中 $G(k,\chi)=\sum_{a=1}^{m-1}\chi(a)\zeta_m^{ak}$ 为 Gauss 和.
\end{lemma}

\begin{proof}
记 $\omega=\zeta_m$, 当 $\operatorname{Re}(s)>1$ 时
\begin{align*}
  L(s,\chi)&=\sum_{a=0}^{m-1}\chi(a)\sum_{\substack{n=1\\ n\equiv a\ (\mathrm{mod}\ m)}}^{\infty}n^{-s}\\
  &=\sum_{a=0}^{m-1}\chi(a)\sum_{n=1}^{\infty}\left(\frac{1}{m}\sum_{k=0}^{m-1}\omega^{(a-n)k}\right)n^{-s}\\
  &=\frac{1}{m}\sum_{k=0}^{m-1}G(k,\chi)\sum_{n=1}^{\infty}\omega^{-nk}n^{-s}\\
  &=\frac{G(1,\chi)}{m}\sum_{k=0}^{m-1}\overline{\chi}(k)\sum_{n=1}^{\infty}\omega^{-nk}n^{-s}.
\end{align*}
（交换和号是由于级数的绝对收敛性）. 但是 D-级数 $\sum_{n=1}^{\infty}\omega^{-nk}n^{-s}$ 在 $\operatorname{Re}(s)>0$ 中正则, 从而它在 $s=1$ 处连续. 熟知它在 $s=1$ 处的值为 $-\log(1-\omega^{-k})$. 因此
\begin{align*}
  L(1,\chi)&=-\frac{G(1,\chi)}{m}\sum_{k=0}^{m-1}\overline{\chi}(k)\log(1-\omega^{-k})\\
  &=-\frac{\chi(-1)G(1,\chi)}{m}\sum_{k=0}^{m-1}\overline{\chi}(k)\log(1-\omega^k).
\end{align*}
当 $0<k<m$ 时,
\[
  \log(1-\omega^k)=\log 2+\log\sin\frac{k\pi}{m}+\left(\frac{k}{m}-\frac{1}{2}\right)\pi i,
\]
于是（由于 $\operatorname{cond}(\chi)=m\geqslant 3$ 从而 $\chi\neq\chi_0$）
\[
  L(1,\chi)=-\frac{\chi(-1)G(1,\chi)}{m}\sum_{k=1}^{m-1}\overline{\chi}(k)\left(\log\sin\frac{k\pi}{m}+\frac{k\pi i}{m}\right).
\]
当 $\chi(-1)=1$ 时, $\sum\chi(k)k=0$, 而当 $\chi(-1)=-1$ 时,
\[
  \sum\overline{\chi}(k)\log\sin\frac{k\pi}{m}=0
\]
（习题）, 从而
\[
  L(1,\chi)=\begin{cases}
    -\dfrac{2G(1,\chi)}{m}\sum\limits_{1\leqslant k<m/2}\overline{\chi}(k)\log\sin\dfrac{k\pi}{m}, & \chi(-1)=1\ \text{时};\\[1em]
    \dfrac{\pi iG(1,\chi)}{m^2}\sum\limits_{k=1}^{m-1}\overline{\chi}(k)k, & \chi(-1)=-1\ \text{时}.
  \end{cases}
\]
最后再由下面引理 \ref{lem:ch6-5}, 即得到引理 \ref{lem:ch6-4} 的全部结果.
\end{proof}

\begin{lemma}\label{lem:ch6-5}
设 $\chi$ 为模 $m$ 的本原奇特征, 则
\[
  \sum_{k=1}^{m-1}\chi(k)k=\frac{m}{\chi(2)-2}\sum_{1\leqslant k<m/2}\chi(k).
\]
\end{lemma}

\begin{proof}
如果 $2\nmid m$, 则
\begin{align*}
  -\sum_{0<k<m}\chi(k)k&=-\sum_{\substack{0<k<m\\ 2\mid k}}\left(\chi(k)k+\chi(m-k)(m-k)\right)\\
  &=-\sum_{0<k<m/2}\chi(2k)2k+\sum_{0<k<m/2}\chi(2k)(m-2k)\\
  &=\chi(2)\sum_{0<k<m/2}(m-4k)\chi(k).
\end{align*}
因此
\begin{equation}\tag{1}
  -\overline{\chi}(2)\sum_{0<k<m}\chi(k)k=\sum_{0<k<m/2}(m-4k)\chi(k)
\end{equation}
另一方面,
\begin{align*}
  -\sum_{0<k<m}\chi(k)k&=-\sum_{0<k<m/2}\chi(k)k-\sum_{0<k<m/2}\chi(m-k)(m-k)\\
  &=-\sum_{0<k<m/2}\chi(k)k+\sum_{0<k<m/2}\chi(k)(m-k)\\
  &=\sum_{0<k<m/2}(m-2k)\chi(k).
\end{align*}
将此式与 (1) 式合在一起, 即为
\[
  \sum_{0<k<m}\chi(k)k=\frac{m}{\chi(2)-2}\sum_{0<k<m/2}\chi(k).
\]
如果 $4\mid m$, 则 $\chi(2)=0$, 并且 $\chi(k+m/2)=-\chi(k)(0<k<m/2)$（这是因为: 由 $(m/2+1)^2\equiv\frac{m}{4}\cdot m+m+1\equiv 1\ (\mathrm{mod}\ m)$ 可知 $\chi(m/2+1)=\pm 1$. 如果 $\chi(m/2+1)=1$, 则由于对每个奇数 $k$ 均有 $(m/2+1)k\equiv m/2+k\ (\mathrm{mod}\ m)$, 从而 $\chi(m/2+k)=\chi(k)$. 于是导出矛盾 $m=\operatorname{cond}(\chi)\mid m/2$. 因此 $\chi(m/2+1)=-1$, 从而对于每个奇数 $k$ 均有 $\chi(k+m/2)=-\chi(k)$. 而对偶数 $k$ 此式当然也是对的, 因两边均为 $0$）. 于是
\begin{align*}
  \sum_{0<k<m}\chi(k)k&=\sum_{0<k<m/2}\chi(k)k+\sum_{0<k<m/2}\chi(k+m/2)(k+m/2)\\
  &=\sum_{0<k<m/2}\chi(k)k-\sum_{0<k<m/2}\chi(k)(k+m/2)\\
  &=-\frac{m}{2}\sum_{0<k<m/2}\chi(k).
\end{align*}
即引理对 $4\mid m$ 情形也成立.
\end{proof}

将引理 \ref{lem:ch6-4} 代入引理 \ref{lem:ch6-3} 就得到

\begin{lemma}\label{lem:ch6-6}
以 $f_\chi$ 表示 D-特征 $\chi$ 的导子, 则

(a) 对于 $n$ 次实 Abel 数域 $K$,
\[
  R_Kh_K=|d(K)|^{1/2}\prod_{\chi_0\neq\chi\in\widehat{K}}\frac{1}{\sqrt{f_\chi}}\left|\sum_{1\leqslant k<f_\chi/2}\chi(k)\log\sin\frac{k\pi}{f_\chi}\right|.
\]
(b) 对于 $n$ 次虚 Abel 数域 $K$,
\[
  R_{K_+}h_{K_+}=|d(K_+)|^{1/2}\prod_{\substack{\chi_0\neq\chi\in\widehat{K}\\ \chi(-1)=1}}\frac{1}{\sqrt{f_\chi}}\left|\sum_{1\leqslant k<f_\chi/2}\chi(k)\log\sin\frac{k\pi}{f_\chi}\right|,
\]
\[
  \begin{aligned}
    \frac{R_Kh_K}{R_{K_+}h_{K_+}}&=\frac{w_K}{2}|d(K)/d(K_+)|^{1/2}\prod_{\substack{\chi\in\widehat{K}\\ \chi(-1)=-1}}\frac{1}{f_\chi^{3/2}}\left|\sum_{k=1}^{f_\chi-1}\chi(k)k\right|\\
    &=\frac{w_K}{2}|d(K)/d(K_+)|^{1/2}\prod_{\substack{\chi\in\widehat{K}\\ \chi(-1)=-1}}\frac{1}{\sqrt{f_\chi}|\chi(2)-2|}\left|\sum_{1\leqslant k<f_\chi/2}\chi(k)\right|.
  \end{aligned}
\]
\end{lemma}

我们仍然可以将上面诸式右边作进一步化简. 为此, 我们需要第四章所介绍的 $L(s,\chi)$ 和 $\zeta_K(s)$ 的函数方程. 它们可导出 $d(K)$ 和 $f_\chi$ 之间一个简单而奇妙的关系, 就是著名的 Hasse 判别式——导子公式.

\begin{lemma}[\emph{Hasse, 判别式——导子公式}]\label{lem:ch6-7}
设 $K$ 为 Abel 数域, 则
\[
  \prod_{\chi\in\widehat{K}}f_\chi=|d(K)|.
\]
\end{lemma}

\begin{proof}
将公式 $\zeta_K(s)=\prod_{\chi\in\widehat{K}}L(s,\chi^*)$ 两边的 D-级数均应用各自的函数方程. 由 $\zeta_K(s)$ 的函数方程不难算出, 若令
\[
  \Phi_K(s)=\begin{cases}
    |d(K)|^{s/2}(\pi^{-s/2}\Gamma(s/2))^n\zeta_K(s), & \text{若 } K\ \text{为实域};\\[0.5em]
    |d(K)|^{s/2}(\pi^{-s/2}\Gamma(s/2))^{n/2}\left(\pi^{-s/2}\Gamma\left(\frac{1+s}{2}\right)\right)^{n/2}\zeta_K(s), & \text{若 } K\ \text{为虚域},
  \end{cases}
\]
则 $|\Phi_K(s)|=|\Phi_K(1-s)|$. 这里我们使用了公式
\[
  \Gamma(s)=\Gamma(s/2)\Gamma\left(\frac{1+s}{2}\right)\cdot 2^{s-1}/\sqrt{\pi}.
\]
另一方面, 由 $L$-函数的函数方程可知, 若令 $\chi$ 为模 $f_\chi$ 的本原特征, 而
\[
  \Phi(s,\chi)=\begin{cases}
    f_\chi^{s/2}(\pi^{-s/2}\Gamma(s/2))L(s,\chi), & \text{若 } \chi(-1)=1;\\[0.5em]
    f_\chi^{s/2}\left(\pi^{-s/2}\Gamma\left(\frac{s+1}{2}\right)\right)L(s,\chi), & \text{若 } \chi(-1)=-1,
  \end{cases}
\]
则 $|\Phi(s,\chi)|=|\Phi(1-s,\chi)|$. 现在由 $\zeta_K(s)=\prod_{\chi\in\widehat{K}}L(s,\chi^*)$ 即知
\[
  \left|\Phi_K(s)\Big/\prod_{\chi\in\widehat{K}}\Phi(s,\chi)\right|=\left|d(K)\Big/\prod_\chi f_\chi\right|^{s/2}.
\]
但是左边将 $s$ 改为 $1-s$ 时是不变的, 于是右边也应如此. 但这只有 $|d(K)|=\prod f_\chi$ 的时候才可能.
\end{proof}

最后, 将引理 \ref{lem:ch6-7} 和引理 \ref{lem:ch6-6} 放在一起, 我们就给出 Hasse 于 $20$ 世纪 $30$ 年代建立的

\begin{theorem}[\emph{Hasse, Abel 域类数解析公式}]\label{thm:6.2}
(a) 若 $K$ 为 $n$ 次实 Abel 域, 则
\[
  R_Kh_K=\prod_{\chi_0\neq\chi\in\widehat{K}}\left|\sum_{1\leqslant k<f_\chi/2}\chi(k)\log\sin\frac{k\pi}{f_\chi}\right|.
\]
(b) 若 $K$ 为 $n$ 次虚 Abel 域, 则
\[
  R_{K_+}h_{K_+}=\prod_{\substack{\chi_0\neq\chi\in\widehat{K}\\ \chi(-1)=1}}\left|\sum_{1\leqslant k<f_\chi/2}\chi(k)\log\sin\frac{k\pi}{f_\chi}\right|,
\]
\[
  \begin{aligned}
    \frac{R_Kh_K}{R_{K_+}h_{K_+}}&=\frac{w_K}{2}\prod_{\substack{\chi\in\widehat{K}\\ \chi(-1)=-1}}\frac{1}{f_\chi}\left|\sum_{k=1}^{f_\chi-1}\chi(k)k\right|\\
    &=\frac{w_K}{2}\prod_{\substack{\chi\in\widehat{K}\\ \chi(-1)=-1}}\frac{1}{|\chi(2)-2|}\left|\sum_{1\leqslant k<f_\chi/2}\chi(k)\right|.
  \end{aligned}
\]
\end{theorem}

在以下两节里, 我们要对于最早研究类数问题的两种数域——二次域和分圆域给出更简单的类数公式和进一步的结果与猜想.
